%*******************************************************
% Abstract
%*******************************************************
%\renewcommand{\abstractname}{Abstract}
\pdfbookmark[1]{Abstract}{Abstract}
% \addcontentsline{toc}{chapter}{\tocEntry{Abstract}}
\begingroup
\let\clearpage\relax
\let\cleardoublepage\relax
\let\cleardoublepage\relax

\chapter*{Resumen}

Las aplicaciones de realidad virtual (VR), aumentada (AR) y sus derivadas,
englobadas dentro del término XR, necesitan métodos para localizar y entender
los movimientos realizados por el usuario y así poder actualizar la simulación
que ejecutan de manera acorde. El problema de localización ha sido uno de los
mayores desafíos a la hora de producir dispositivos de realidad virtual para
electrónica de consumo en los últimos años. Afortunadamente, han ocurrido
grandes avances en la localización visual-inercial mediante cámaras y sensores
inerciales; dados principalmente gracias a su continua mejora y abaratamiento
relacionados con su producción masiva para telefonía móvil. A pesar de esto,
sistemas de este tipo que puedan ser utilizados para XR se desarrollan
mayormente como parte de productos privativos. En este trabajo, se estudian
sistemas de localización visual-inercial provenientes de la academia y se
detallan las dificultades propias de integrarlos y adaptarlos para aplicaciones
de XR. En particular se los integra con Monado, la implementación de código
libre del estándar OpenXR. De esta forma, se presenta una de las primeras
soluciones completamente libres que permite localizar dispositivos XR con este
tipo de métodos visual-inerciales en sistemas operativos basados en GNU/Linux.

\vfill

\begin{otherlanguage}{american}
\chapter*{Abstract}

Applications for virtual reality, augmented reality, and their
derivatives, all of them encompassed in the term XR, need ways to keep track of
motion in the physical world to be able to reflect these movements accordingly
in the simulated world. Tracking has been one of the main
challenges in bringing VR devices into mainstream consumer electronics in recent
years. Fortunately, there has been breakthroughs in visual-inertial
tracking with cameras and inertial sensors due to their continuous improvement and
reduced costs thanks to the massive production tied to the mobile phone industry. However,
systems that perform this kind of tracking are developed for mostly privative
products. In this work, we study visual-inertial tracking systems that come
from academia, and detail the difficulties that come from integrating and
adapting them for XR applications. In particular, we integrate them to Monado,
the open-source implementation of the OpenXR standard. Therefore, we present one
of the first completely open solutions that integrates this kind of tracking for
XR applications and devices on GNU/Linux-based operating systems.

\end{otherlanguage}

\endgroup

\vfill
